\documentclass[10pt]{article}
\usepackage[a4paper,margin=1.5cm]{geometry}\usepackage{graphicx,amsmath,amssymb,xcolor,hyperref,booktabs,longtable}
\usepackage{fontspec}
\setlength{\parindent}{0pt}
\begin{document}
{\LARGE\bfseries Alternative storyboard v2: ``Where two mirrors agree''}\\[3pt]
Reply Short to Inigo Quilez: \url{https://www.youtube.com/shorts/cc-Z_igk4Lo}. Polished version of the alternative story, with a new first-person opening. \texttt{alt\_storyboard.pdf} and v4 are unchanged.\\[6pt]
\textbf{Format.} YouTube Short, 9:16 (1080$\times$1920). \textbf{174\,s} (cap 179\,s), 39 scenes, 434 words of voiceover ($\approx$2.5 words/s). No audio or MP4 at this stage.\\
\textbf{Look.} Same chalkboard as v4. Mirrors are chalk-white, reflected circles \textcolor[HTML]{4DA3FF}{\textbf{blue}} (dashed = image), \textcolor[HTML]{FF9A3C}{\textbf{orange = imaginary / hidden}}, grey = circles fixed by both mirrors. Dots = crossings, squares = limiting points.\\
\textbf{Opening.} Spoken in first person by Carlo, warm and respectful, with the credit caption ``Reply to Inigo Quilez --- youtube.com/shorts/cc-Z\_igk4Lo'' on screen.\\ \textbf{Through-line.} A circle, a mirror-circle, another circle: a closed algebra (a kei), with an invariant (the angle $\cos^2\theta$, preserved by every mirror). The hidden points are where two mirrors agree, and that agreement set moves with any mirror (the invariant answer). Two mirrors always agree at exactly two real points: the crossings when they cut, the limiting points when they miss. Over $\mathbb C$ there are four, and exactly two are ever real.\\[4pt]
\textbf{Guardrails.} The operation is $C\star D=D-2\langle D,C\rangle C/Q(C)$ on $[a:b:c:d]\in\mathbb P^3$, defined only for $Q(C)\neq0$, so points cannot be mirrors. $C\star C=C$ holds projectively. ``Always two'' means distinct real mirrors, counting $\infty$ (two lines agree at $\infty$), and one point if they are tangent. ``Agree'' means $K\star P=L\star P$, which is the same as $P$ being a fixed point of $s_L\circ s_K$. The invariant is the angle, not $Q$ itself: a circle's radius changes under inversion, and only the vector $Q$ and its sign survive. All frames are in \texttt{alt\_frames\_v2/}, made by \texttt{alt\_frames\_v2.py}; every new number is checked in \texttt{alt\_frames\_v2/check\_log.txt}. Theory: \texttt{alt\_quandle\_note.pdf}.\\[4pt]
{\small\begin{tabular}{lll}\toprule act & scenes & time\\\midrule Opening (first person) & 1--3 & 0:00--0:12 \\ His move, my move & 4--5 & 0:12--0:21 \\ Circles as mirrors & 6--12 & 0:21--0:50 \\ The algebra: a kei & 13--17 & 0:50--1:13 \\ The invariant & 18--19 & 1:13--1:23 \\ What a mirror fixes & 20--21 & 1:23--1:31 \\ Two mirrors agree & 22--33 & 1:31--2:27 \\ Four, two real; the invariant answer; close & 34--39 & 2:27--2:54 \\\bottomrule\end{tabular}}\\[4pt]
{\scriptsize\begin{longtable}{rllp{0.55\textwidth}}\toprule \# & time & beat & voiceover\\\midrule
1 & 0:00--0:04 & A lovely video & \emph{Inigo Quilez's video is lovely, and it opens many doors.}\\
2 & 0:04--0:08 & Circles as mirrors & \emph{I took one: circles as mirrors, the Möbius way.}\\
3 & 0:08--0:12 & An algebra, and an invariant & \emph{My aim: give his hidden points an algebra, and an invariant.}\\
4 & 0:12--0:17 & His move & \emph{His move: when the square root goes negative, go complex. Two orange dots.}\\
5 & 0:17--0:21 & My move & \emph{My move: a circle, a mirror-circle, and out comes another circle.}\\
6 & 0:21--0:25 & A line is a mirror & \emph{A line is the simplest mirror: reflect a circle across it.}\\
7 & 0:25--0:29 & A circle is a mirror too & \emph{A circle is a mirror too. That reflection is called inversion.}\\
8 & 0:29--0:33 & Line in circle & \emph{Invert the line in the circle: out comes a small circle.}\\
9 & 0:33--0:38 & Order matters & \emph{Reflect the circle in the line instead: a different circle. Order matters.}\\
10 & 0:38--0:42 & Four numbers & \emph{Every circle is just four numbers: a, b, c, d.}\\
11 & 0:42--0:46 & One formula & \emph{And reflection is a single formula on those numbers.}\\
12 & 0:46--0:50 & Closed & \emph{Circle in, circle out. Lines too. The family is closed.}\\
13 & 0:50--0:55 & Rules 1 and 2 & \emph{Rule one: a mirror fixes itself. Rule two: reflect twice, and you're back.}\\
14 & 0:55--1:00 & Rule 3 & \emph{Rule three: reflect a whole picture, and every relation in it comes along.}\\
15 & 1:00--1:04 & A kei & \emph{These rules have a name: a kei, or involutory quandle.}\\
16 & 1:04--1:08 & Not associative & \emph{Not associative, though: move the brackets, and the circle changes.}\\
17 & 1:08--1:13 & Points & \emph{Points can't be mirrors, since everything would collapse. But they still get reflected.}\\
18 & 1:13--1:18 & The invariant & \emph{Now, the invariant: every mirror keeps the angle between any two circles.}\\
19 & 1:18--1:23 & Hidden stays hidden & \emph{And that number decides real or hidden, so no mirror can flip it.}\\
20 & 1:23--1:27 & What a mirror fixes & \emph{What does a mirror leave alone? The points on it.}\\
21 & 1:27--1:31 & ...and right angles & \emph{And every circle that crosses it at right angles.}\\
22 & 1:31--1:35 & Two mirrors & \emph{Now take two mirrors: our circle and our line.}\\
23 & 1:35--1:39 & Where do they agree? & \emph{Where do they agree? Same point in, same point out.}\\
24 & 1:39--1:44 & They agree at the crossings & \emph{When the line cuts, they agree at the crossings. Both leave them fixed.}\\
25 & 1:44--1:48 & Fixed by both & \emph{The circles both mirrors fix swirl around those two points.}\\
26 & 1:48--1:53 & Lift the line & \emph{Lift the line, and the crossings hide. Do the mirrors still agree?}\\
27 & 1:53--1:58 & They still agree & \emph{They do, at two real points on the axis: the limiting points.}\\
28 & 1:58--2:03 & Both swap them & \emph{Each mirror swaps them, top for bottom, and their heights multiply to one.}\\
29 & 2:03--2:07 & Fixed by both, again & \emph{Now the circles both mirrors fix all run through them.}\\
30 & 2:07--2:11 & Quilez's dots & \emph{And Quilez's dots? The same two points, turned ninety degrees.}\\
31 & 2:11--2:16 & Sweep the line & \emph{Sweep the line: the points slide up, merge where it touches, then split.}\\
32 & 2:16--2:21 & Always two & \emph{Two mirrors always agree at exactly two points, or one if they touch.}\\
33 & 2:21--2:27 & Why two? & \emph{Why two? Two reflections make one Möbius map. Its fixed points solve a quadratic.}\\
34 & 2:27--2:32 & Over C: four & \emph{Over the complex numbers there are four: two crossings, two limiting points.}\\
35 & 2:32--2:36 & Exactly two are real & \emph{But only two are ever real. The other two hide.}\\
36 & 2:36--2:41 & The invariant answer & \emph{Move everything by any mirror, and agreement points follow. That's the invariant answer.}\\
37 & 2:41--2:46 & Partners & \emph{So the hidden pair isn't lost. It's the hidden half of the agreement.}\\
38 & 2:46--2:50 & A closed algebra & \emph{A circle, a mirror-circle, another circle: one closed algebra.}\\
39 & 2:50--2:54 & Where two mirrors agree & \emph{Not hidden. Just where two mirrors agree.}\\\bottomrule\end{longtable}}
\newpage
\begin{minipage}[t]{0.17\textwidth}\vspace{0pt}\includegraphics[width=\linewidth]{alt_frames_v2/s01.png}\end{minipage}\hfill
\begin{minipage}[t]{0.80\textwidth}\vspace{0pt}\small
{\bfseries Scene 1: A lovely video}\hfill 0:00--0:04 (4\,s)\\[2pt]
\textbf{VO:} \emph{Inigo Quilez's video is lovely, and it opens many doors.}\\[2pt]
\textbf{On screen:} ``Reply to Inigo Quilez --- youtube.com/shorts/cc-Z\_igk4Lo''\\[2pt]
\textbf{Animation:} His setup, redrawn: unit circle, a tilted line that misses it, two orange hidden dots. Credit caption fades in. Spoken in first person by Carlo.
\end{minipage}\par\vspace{4pt}\hrule\vspace{4pt}
\begin{minipage}[t]{0.17\textwidth}\vspace{0pt}\includegraphics[width=\linewidth]{alt_frames_v2/s02.png}\end{minipage}\hfill
\begin{minipage}[t]{0.80\textwidth}\vspace{0pt}\small
{\bfseries Scene 2: Circles as mirrors}\hfill 0:04--0:08 (4\,s)\\[2pt]
\textbf{VO:} \emph{I took one: circles as mirrors, the Möbius way.}\\[2pt]
\textbf{On screen:} ``the Möbius way''\\[2pt]
\textbf{Animation:} The same circle and line turn chalk-white (mirrors). The circle is reflected in the line, and the line is inverted in the circle (dashed).
\end{minipage}\par\vspace{4pt}\hrule\vspace{4pt}
\begin{minipage}[t]{0.17\textwidth}\vspace{0pt}\includegraphics[width=\linewidth]{alt_frames_v2/s03.png}\end{minipage}\hfill
\begin{minipage}[t]{0.80\textwidth}\vspace{0pt}\small
{\bfseries Scene 3: An algebra, and an invariant}\hfill 0:08--0:12 (4\,s)\\[2pt]
\textbf{VO:} \emph{My aim: give his hidden points an algebra, and an invariant.}\\[2pt]
\textbf{On screen:} $C\star D$ $\cdot$ $\cos^2\theta$\\[2pt]
\textbf{Animation:} The formula and the invariant appear under the picture. The orange dots stay.
\end{minipage}\par\vspace{4pt}\hrule\vspace{4pt}
\begin{minipage}[t]{0.17\textwidth}\vspace{0pt}\includegraphics[width=\linewidth]{alt_frames_v2/s04.png}\end{minipage}\hfill
\begin{minipage}[t]{0.80\textwidth}\vspace{0pt}\small
{\bfseries Scene 4: His move}\hfill 0:12--0:17 (5\,s)\\[2pt]
\textbf{VO:} \emph{His move: when the square root goes negative, go complex. Two orange dots.}\\[2pt]
\textbf{On screen:} $\pm i\sqrt{|\Delta|}$\\[2pt]
\textbf{Animation:} Reused v4 frame: $\pm i$ flashes; two orange dots drop onto the line.
\end{minipage}\par\vspace{4pt}\hrule\vspace{4pt}
\begin{minipage}[t]{0.17\textwidth}\vspace{0pt}\includegraphics[width=\linewidth]{alt_frames_v2/s05.png}\end{minipage}\hfill
\begin{minipage}[t]{0.80\textwidth}\vspace{0pt}\small
{\bfseries Scene 5: My move}\hfill 0:17--0:21 (4\,s)\\[2pt]
\textbf{VO:} \emph{My move: a circle, a mirror-circle, and out comes another circle.}\\[2pt]
\textbf{On screen:} $K\star D$\\[2pt]
\textbf{Animation:} Blue circle $D$ flips through $K$ into the dashed image.
\end{minipage}\par\vspace{4pt}\hrule\vspace{4pt}
\begin{minipage}[t]{0.17\textwidth}\vspace{0pt}\includegraphics[width=\linewidth]{alt_frames_v2/s06.png}\end{minipage}\hfill
\begin{minipage}[t]{0.80\textwidth}\vspace{0pt}\small
{\bfseries Scene 6: A line is a mirror}\hfill 0:21--0:25 (4\,s)\\[2pt]
\textbf{VO:} \emph{A line is the simplest mirror: reflect a circle across it.}\\[2pt]
\textbf{On screen:} $L\star D$\\[2pt]
\textbf{Animation:} Circle at $(0.9,0.6)$ reflects across $y=1.7$ to $(0.9,2.8)$.
\end{minipage}\par\vspace{4pt}\hrule\vspace{4pt}
\begin{minipage}[t]{0.17\textwidth}\vspace{0pt}\includegraphics[width=\linewidth]{alt_frames_v2/s07.png}\end{minipage}\hfill
\begin{minipage}[t]{0.80\textwidth}\vspace{0pt}\small
{\bfseries Scene 7: A circle is a mirror too}\hfill 0:25--0:29 (4\,s)\\[2pt]
\textbf{VO:} \emph{A circle is a mirror too. That reflection is called inversion.}\\[2pt]
\textbf{On screen:} $p\mapsto p/|p|^2$\\[2pt]
\textbf{Animation:} Rays from the centre; each point jumps to its partner on the same ray.
\end{minipage}\par\vspace{4pt}\hrule\vspace{4pt}
\begin{minipage}[t]{0.17\textwidth}\vspace{0pt}\includegraphics[width=\linewidth]{alt_frames_v2/s08.png}\end{minipage}\hfill
\begin{minipage}[t]{0.80\textwidth}\vspace{0pt}\small
{\bfseries Scene 8: Line in circle}\hfill 0:29--0:33 (4\,s)\\[2pt]
\textbf{VO:} \emph{Invert the line in the circle: out comes a small circle.}\\[2pt]
\textbf{On screen:} $K\star L$: centre $(0,0.294)$, $r=1/3.4$\\[2pt]
\textbf{Animation:} The line curls into a small circle through the centre.
\end{minipage}\par\vspace{4pt}\hrule\vspace{4pt}
\begin{minipage}[t]{0.17\textwidth}\vspace{0pt}\includegraphics[width=\linewidth]{alt_frames_v2/s09.png}\end{minipage}\hfill
\begin{minipage}[t]{0.80\textwidth}\vspace{0pt}\small
{\bfseries Scene 9: Order matters}\hfill 0:33--0:38 (5\,s)\\[2pt]
\textbf{VO:} \emph{Reflect the circle in the line instead: a different circle. Order matters.}\\[2pt]
\textbf{On screen:} $K\star L\neq L\star K$\\[2pt]
\textbf{Animation:} $L\star K$ (centre $(0,3.4)$, $r=1$) appears beside $K\star L$.
\end{minipage}\par\vspace{4pt}\hrule\vspace{4pt}
\begin{minipage}[t]{0.17\textwidth}\vspace{0pt}\includegraphics[width=\linewidth]{alt_frames_v2/s10.png}\end{minipage}\hfill
\begin{minipage}[t]{0.80\textwidth}\vspace{0pt}\small
{\bfseries Scene 10: Four numbers}\hfill 0:38--0:42 (4\,s)\\[2pt]
\textbf{VO:} \emph{Every circle is just four numbers: a, b, c, d.}\\[2pt]
\textbf{On screen:} $[a:b:c:d]$\\[2pt]
\textbf{Animation:} Reused v4 frame: coefficients lift into a bracket.
\end{minipage}\par\vspace{4pt}\hrule\vspace{4pt}
\begin{minipage}[t]{0.17\textwidth}\vspace{0pt}\includegraphics[width=\linewidth]{alt_frames_v2/s11.png}\end{minipage}\hfill
\begin{minipage}[t]{0.80\textwidth}\vspace{0pt}\small
{\bfseries Scene 11: One formula}\hfill 0:42--0:46 (4\,s)\\[2pt]
\textbf{VO:} \emph{And reflection is a single formula on those numbers.}\\[2pt]
\textbf{On screen:} $C\star D=D-\frac{2\langle D,C\rangle}{Q(C)}C$\\[2pt]
\textbf{Animation:} Formula slams in over the scene-5 picture.
\end{minipage}\par\vspace{4pt}\hrule\vspace{4pt}
\begin{minipage}[t]{0.17\textwidth}\vspace{0pt}\includegraphics[width=\linewidth]{alt_frames_v2/s12.png}\end{minipage}\hfill
\begin{minipage}[t]{0.80\textwidth}\vspace{0pt}\small
{\bfseries Scene 12: Closed}\hfill 0:46--0:50 (4\,s)\\[2pt]
\textbf{VO:} \emph{Circle in, circle out. Lines too. The family is closed.}\\[2pt]
\textbf{On screen:} circle $\star$ circle $=$ circle\\[2pt]
\textbf{Animation:} Nine random circles and a line inverted at once.
\end{minipage}\par\vspace{4pt}\hrule\vspace{4pt}
\begin{minipage}[t]{0.17\textwidth}\vspace{0pt}\includegraphics[width=\linewidth]{alt_frames_v2/s13.png}\end{minipage}\hfill
\begin{minipage}[t]{0.80\textwidth}\vspace{0pt}\small
{\bfseries Scene 13: Rules 1 and 2}\hfill 0:50--0:55 (5\,s)\\[2pt]
\textbf{VO:} \emph{Rule one: a mirror fixes itself. Rule two: reflect twice, and you're back.}\\[2pt]
\textbf{On screen:} $C\star C=C$; $C\star(C\star D)=D$\\[2pt]
\textbf{Animation:} Left: $K$ reflected in itself lands on itself. Right: out-and-back arrows.
\end{minipage}\par\vspace{4pt}\hrule\vspace{4pt}
\begin{minipage}[t]{0.17\textwidth}\vspace{0pt}\includegraphics[width=\linewidth]{alt_frames_v2/s14.png}\end{minipage}\hfill
\begin{minipage}[t]{0.80\textwidth}\vspace{0pt}\small
{\bfseries Scene 14: Rule 3}\hfill 0:55--1:00 (5\,s)\\[2pt]
\textbf{VO:} \emph{Rule three: reflect a whole picture, and every relation in it comes along.}\\[2pt]
\textbf{On screen:} $C\star(D\star E)=(C\star D)\star(C\star E)$\\[2pt]
\textbf{Animation:} The picture $D,E,D\star E$ flips across $y=1$ intact.
\end{minipage}\par\vspace{4pt}\hrule\vspace{4pt}
\begin{minipage}[t]{0.17\textwidth}\vspace{0pt}\includegraphics[width=\linewidth]{alt_frames_v2/s15.png}\end{minipage}\hfill
\begin{minipage}[t]{0.80\textwidth}\vspace{0pt}\small
{\bfseries Scene 15: A kei}\hfill 1:00--1:04 (4\,s)\\[2pt]
\textbf{VO:} \emph{These rules have a name: a kei, or involutory quandle.}\\[2pt]
\textbf{On screen:} three axioms; Takasaki 1943, Joyce 1982\\[2pt]
\textbf{Animation:} Axioms type on one per beat.
\end{minipage}\par\vspace{4pt}\hrule\vspace{4pt}
\begin{minipage}[t]{0.17\textwidth}\vspace{0pt}\includegraphics[width=\linewidth]{alt_frames_v2/s16.png}\end{minipage}\hfill
\begin{minipage}[t]{0.80\textwidth}\vspace{0pt}\small
{\bfseries Scene 16: Not associative}\hfill 1:04--1:08 (4\,s)\\[2pt]
\textbf{VO:} \emph{Not associative, though: move the brackets, and the circle changes.}\\[2pt]
\textbf{On screen:} radii $0.057$ vs $0.125$\\[2pt]
\textbf{Animation:} Inset zoom shows the two results.
\end{minipage}\par\vspace{4pt}\hrule\vspace{4pt}
\begin{minipage}[t]{0.17\textwidth}\vspace{0pt}\includegraphics[width=\linewidth]{alt_frames_v2/s17.png}\end{minipage}\hfill
\begin{minipage}[t]{0.80\textwidth}\vspace{0pt}\small
{\bfseries Scene 17: Points}\hfill 1:08--1:13 (5\,s)\\[2pt]
\textbf{VO:} \emph{Points can't be mirrors, since everything would collapse. But they still get reflected.}\\[2pt]
\textbf{On screen:} $Q(C)=0$: no mirror; $K\star P$ is a point\\[2pt]
\textbf{Animation:} Left: the image collapses as the mirror shrinks. Right: points jump to their inverse points.
\end{minipage}\par\vspace{4pt}\hrule\vspace{4pt}
\begin{minipage}[t]{0.17\textwidth}\vspace{0pt}\includegraphics[width=\linewidth]{alt_frames_v2/s18.png}\end{minipage}\hfill
\begin{minipage}[t]{0.80\textwidth}\vspace{0pt}\small
{\bfseries Scene 18: The invariant}\hfill 1:13--1:18 (5\,s)\\[2pt]
\textbf{VO:} \emph{Now, the invariant: every mirror keeps the angle between any two circles.}\\[2pt]
\textbf{On screen:} $I=\frac{\langle X,Y\rangle^2}{Q(X)Q(Y)}=\cos^2\theta$: $0.0156\to0.0156$\\[2pt]
\textbf{Animation:} Two crossing circles ($\cos\theta=0.125$) and their images in a mirror circle; the angle is marked on both.
\end{minipage}\par\vspace{4pt}\hrule\vspace{4pt}
\begin{minipage}[t]{0.17\textwidth}\vspace{0pt}\includegraphics[width=\linewidth]{alt_frames_v2/s19.png}\end{minipage}\hfill
\begin{minipage}[t]{0.80\textwidth}\vspace{0pt}\small
{\bfseries Scene 19: Hidden stays hidden}\hfill 1:18--1:23 (5\,s)\\[2pt]
\textbf{VO:} \emph{And that number decides real or hidden, so no mirror can flip it.}\\[2pt]
\textbf{On screen:} $I=2.89>1$ before and after\\[2pt]
\textbf{Animation:} The circle, the line $y=1.7$ and the orange hidden pair, with their images in a mirror. The images still miss, and the pair is still orange.
\end{minipage}\par\vspace{4pt}\hrule\vspace{4pt}
\begin{minipage}[t]{0.17\textwidth}\vspace{0pt}\includegraphics[width=\linewidth]{alt_frames_v2/s20.png}\end{minipage}\hfill
\begin{minipage}[t]{0.80\textwidth}\vspace{0pt}\small
{\bfseries Scene 20: What a mirror fixes}\hfill 1:23--1:27 (4\,s)\\[2pt]
\textbf{VO:} \emph{What does a mirror leave alone? The points on it.}\\[2pt]
\textbf{On screen:} $K\star P=P\Leftrightarrow\langle P,K\rangle=0$\\[2pt]
\textbf{Animation:} Points on $K$ pulse and stay put.
\end{minipage}\par\vspace{4pt}\hrule\vspace{4pt}
\begin{minipage}[t]{0.17\textwidth}\vspace{0pt}\includegraphics[width=\linewidth]{alt_frames_v2/s21.png}\end{minipage}\hfill
\begin{minipage}[t]{0.80\textwidth}\vspace{0pt}\small
{\bfseries Scene 21: ...and right angles}\hfill 1:27--1:31 (4\,s)\\[2pt]
\textbf{VO:} \emph{And every circle that crosses it at right angles.}\\[2pt]
\textbf{On screen:} $\mathrm{Fix}(K)=\{K\}\cup K^\perp$\\[2pt]
\textbf{Animation:} Grey orthogonal circles appear.
\end{minipage}\par\vspace{4pt}\hrule\vspace{4pt}
\begin{minipage}[t]{0.17\textwidth}\vspace{0pt}\includegraphics[width=\linewidth]{alt_frames_v2/s22.png}\end{minipage}\hfill
\begin{minipage}[t]{0.80\textwidth}\vspace{0pt}\small
{\bfseries Scene 22: Two mirrors}\hfill 1:31--1:35 (4\,s)\\[2pt]
\textbf{VO:} \emph{Now take two mirrors: our circle and our line.}\\[2pt]
\textbf{On screen:} $K$, $L$: $y=0.6$\\[2pt]
\textbf{Animation:} Line drops to $c=0.6$.
\end{minipage}\par\vspace{4pt}\hrule\vspace{4pt}
\begin{minipage}[t]{0.17\textwidth}\vspace{0pt}\includegraphics[width=\linewidth]{alt_frames_v2/s23.png}\end{minipage}\hfill
\begin{minipage}[t]{0.80\textwidth}\vspace{0pt}\small
{\bfseries Scene 23: Where do they agree?}\hfill 1:35--1:39 (4\,s)\\[2pt]
\textbf{VO:} \emph{Where do they agree? Same point in, same point out.}\\[2pt]
\textbf{On screen:} $K\star P=L\star P\,?$\\[2pt]
\textbf{Animation:} A test point gives two different images.
\end{minipage}\par\vspace{4pt}\hrule\vspace{4pt}
\begin{minipage}[t]{0.17\textwidth}\vspace{0pt}\includegraphics[width=\linewidth]{alt_frames_v2/s24.png}\end{minipage}\hfill
\begin{minipage}[t]{0.80\textwidth}\vspace{0pt}\small
{\bfseries Scene 24: They agree at the crossings}\hfill 1:39--1:44 (5\,s)\\[2pt]
\textbf{VO:} \emph{When the line cuts, they agree at the crossings. Both leave them fixed.}\\[2pt]
\textbf{On screen:} crossings $(\pm0.8,0.6)$\\[2pt]
\textbf{Animation:} Blue dots pulse; nothing moves.
\end{minipage}\par\vspace{4pt}\hrule\vspace{4pt}
\begin{minipage}[t]{0.17\textwidth}\vspace{0pt}\includegraphics[width=\linewidth]{alt_frames_v2/s25.png}\end{minipage}\hfill
\begin{minipage}[t]{0.80\textwidth}\vspace{0pt}\small
{\bfseries Scene 25: Fixed by both}\hfill 1:44--1:48 (4\,s)\\[2pt]
\textbf{VO:} \emph{The circles both mirrors fix swirl around those two points.}\\[2pt]
\textbf{On screen:} $\mathrm{Fix}(K)\cap\mathrm{Fix}(L)$\\[2pt]
\textbf{Animation:} Grey Apollonian circles.
\end{minipage}\par\vspace{4pt}\hrule\vspace{4pt}
\begin{minipage}[t]{0.17\textwidth}\vspace{0pt}\includegraphics[width=\linewidth]{alt_frames_v2/s26.png}\end{minipage}\hfill
\begin{minipage}[t]{0.80\textwidth}\vspace{0pt}\small
{\bfseries Scene 26: Lift the line}\hfill 1:48--1:53 (5\,s)\\[2pt]
\textbf{VO:} \emph{Lift the line, and the crossings hide. Do the mirrors still agree?}\\[2pt]
\textbf{On screen:} $(\pm1.375i,\,1.7)$\\[2pt]
\textbf{Animation:} Line lifts; blue dots become orange rings.
\end{minipage}\par\vspace{4pt}\hrule\vspace{4pt}
\begin{minipage}[t]{0.17\textwidth}\vspace{0pt}\includegraphics[width=\linewidth]{alt_frames_v2/s27.png}\end{minipage}\hfill
\begin{minipage}[t]{0.80\textwidth}\vspace{0pt}\small
{\bfseries Scene 27: They still agree}\hfill 1:53--1:58 (5\,s)\\[2pt]
\textbf{VO:} \emph{They do, at two real points on the axis: the limiting points.}\\[2pt]
\textbf{On screen:} $(0,3.075)$, $(0,0.325)$\\[2pt]
\textbf{Animation:} Blue squares appear.
\end{minipage}\par\vspace{4pt}\hrule\vspace{4pt}
\begin{minipage}[t]{0.17\textwidth}\vspace{0pt}\includegraphics[width=\linewidth]{alt_frames_v2/s28.png}\end{minipage}\hfill
\begin{minipage}[t]{0.80\textwidth}\vspace{0pt}\small
{\bfseries Scene 28: Both swap them}\hfill 1:58--2:03 (5\,s)\\[2pt]
\textbf{VO:} \emph{Each mirror swaps them, top for bottom, and their heights multiply to one.}\\[2pt]
\textbf{On screen:} $K\star P_+=P_-=L\star P_+$; $3.0748\times0.3252=1$\\[2pt]
\textbf{Animation:} Two arrows swap the squares; the product writes on.
\end{minipage}\par\vspace{4pt}\hrule\vspace{4pt}
\begin{minipage}[t]{0.17\textwidth}\vspace{0pt}\includegraphics[width=\linewidth]{alt_frames_v2/s29.png}\end{minipage}\hfill
\begin{minipage}[t]{0.80\textwidth}\vspace{0pt}\small
{\bfseries Scene 29: Fixed by both, again}\hfill 2:03--2:07 (4\,s)\\[2pt]
\textbf{VO:} \emph{Now the circles both mirrors fix all run through them.}\\[2pt]
\textbf{On screen:} $\mathrm{Fix}(K)\cap\mathrm{Fix}(L)$\\[2pt]
\textbf{Animation:} Grey circles pass through both squares.
\end{minipage}\par\vspace{4pt}\hrule\vspace{4pt}
\begin{minipage}[t]{0.17\textwidth}\vspace{0pt}\includegraphics[width=\linewidth]{alt_frames_v2/s30.png}\end{minipage}\hfill
\begin{minipage}[t]{0.80\textwidth}\vspace{0pt}\small
{\bfseries Scene 30: Quilez's dots}\hfill 2:07--2:11 (4\,s)\\[2pt]
\textbf{VO:} \emph{And Quilez's dots? The same two points, turned ninety degrees.}\\[2pt]
\textbf{On screen:} rotate $90^\circ$ about $(0,c)$\\[2pt]
\textbf{Animation:} Orange dots rotate onto the squares.
\end{minipage}\par\vspace{4pt}\hrule\vspace{4pt}
\begin{minipage}[t]{0.17\textwidth}\vspace{0pt}\includegraphics[width=\linewidth]{alt_frames_v2/s31.png}\end{minipage}\hfill
\begin{minipage}[t]{0.80\textwidth}\vspace{0pt}\small
{\bfseries Scene 31: Sweep the line}\hfill 2:11--2:16 (5\,s)\\[2pt]
\textbf{VO:} \emph{Sweep the line: the points slide up, merge where it touches, then split.}\\[2pt]
\textbf{On screen:} $c:0\to2.6$\\[2pt]
\textbf{Animation:} Trail of agreement points.
\end{minipage}\par\vspace{4pt}\hrule\vspace{4pt}
\begin{minipage}[t]{0.17\textwidth}\vspace{0pt}\includegraphics[width=\linewidth]{alt_frames_v2/s32.png}\end{minipage}\hfill
\begin{minipage}[t]{0.80\textwidth}\vspace{0pt}\small
{\bfseries Scene 32: Always two}\hfill 2:16--2:21 (5\,s)\\[2pt]
\textbf{VO:} \emph{Two mirrors always agree at exactly two points, or one if they touch.}\\[2pt]
\textbf{On screen:} \# $=2$ (1 if tangent); 2000 random pairs, 0 exceptions\\[2pt]
\textbf{Animation:} Two random pairs: crossing, disjoint.
\end{minipage}\par\vspace{4pt}\hrule\vspace{4pt}
\begin{minipage}[t]{0.17\textwidth}\vspace{0pt}\includegraphics[width=\linewidth]{alt_frames_v2/s33.png}\end{minipage}\hfill
\begin{minipage}[t]{0.80\textwidth}\vspace{0pt}\small
{\bfseries Scene 33: Why two?}\hfill 2:21--2:27 (6\,s)\\[2pt]
\textbf{VO:} \emph{Why two? Two reflections make one Möbius map. Its fixed points solve a quadratic.}\\[2pt]
\textbf{On screen:} $z^2-2icz-1=0$\\[2pt]
\textbf{Animation:} Quadratic writes on; roots land on the squares.
\end{minipage}\par\vspace{4pt}\hrule\vspace{4pt}
\begin{minipage}[t]{0.17\textwidth}\vspace{0pt}\includegraphics[width=\linewidth]{alt_frames_v2/s34.png}\end{minipage}\hfill
\begin{minipage}[t]{0.80\textwidth}\vspace{0pt}\small
{\bfseries Scene 34: Over C: four}\hfill 2:27--2:32 (5\,s)\\[2pt]
\textbf{VO:} \emph{Over the complex numbers there are four: two crossings, two limiting points.}\\[2pt]
\textbf{On screen:} $\mathrm{Agree}(K,L)$: 4 points over $\mathbb C$\\[2pt]
\textbf{Animation:} Orange rings and blue squares together.
\end{minipage}\par\vspace{4pt}\hrule\vspace{4pt}
\begin{minipage}[t]{0.17\textwidth}\vspace{0pt}\includegraphics[width=\linewidth]{alt_frames_v2/s35.png}\end{minipage}\hfill
\begin{minipage}[t]{0.80\textwidth}\vspace{0pt}\small
{\bfseries Scene 35: Exactly two are real}\hfill 2:32--2:36 (4\,s)\\[2pt]
\textbf{VO:} \emph{But only two are ever real. The other two hide.}\\[2pt]
\textbf{On screen:} $c=0.6$ | $c=1.7$\\[2pt]
\textbf{Animation:} Split screen: which pair is real swaps.
\end{minipage}\par\vspace{4pt}\hrule\vspace{4pt}
\begin{minipage}[t]{0.17\textwidth}\vspace{0pt}\includegraphics[width=\linewidth]{alt_frames_v2/s36.png}\end{minipage}\hfill
\begin{minipage}[t]{0.80\textwidth}\vspace{0pt}\small
{\bfseries Scene 36: The invariant answer}\hfill 2:36--2:41 (5\,s)\\[2pt]
\textbf{VO:} \emph{Move everything by any mirror, and agreement points follow. That's the invariant answer.}\\[2pt]
\textbf{On screen:} $\mathrm{Agree}(M\star K,M\star L)=M\star\mathrm{Agree}(K,L)$\\[2pt]
\textbf{Animation:} A mirror $M$ (centre $(1.5,0.8)$, $r=1$) carries $K$, $L$ and both squares. The images of the squares are exactly where the new mirrors agree.
\end{minipage}\par\vspace{4pt}\hrule\vspace{4pt}
\begin{minipage}[t]{0.17\textwidth}\vspace{0pt}\includegraphics[width=\linewidth]{alt_frames_v2/s37.png}\end{minipage}\hfill
\begin{minipage}[t]{0.80\textwidth}\vspace{0pt}\small
{\bfseries Scene 37: Partners}\hfill 2:41--2:46 (5\,s)\\[2pt]
\textbf{VO:} \emph{So the hidden pair isn't lost. It's the hidden half of the agreement.}\\[2pt]
\textbf{On screen:} hidden pair $\leftrightarrow$ real agreement pair\\[2pt]
\textbf{Animation:} Dotted links between rings and squares.
\end{minipage}\par\vspace{4pt}\hrule\vspace{4pt}
\begin{minipage}[t]{0.17\textwidth}\vspace{0pt}\includegraphics[width=\linewidth]{alt_frames_v2/s38.png}\end{minipage}\hfill
\begin{minipage}[t]{0.80\textwidth}\vspace{0pt}\small
{\bfseries Scene 38: A closed algebra}\hfill 2:46--2:50 (4\,s)\\[2pt]
\textbf{VO:} \emph{A circle, a mirror-circle, another circle: one closed algebra.}\\[2pt]
\textbf{On screen:} $C\star D$; a kei\\[2pt]
\textbf{Animation:} Recap picture.
\end{minipage}\par\vspace{4pt}\hrule\vspace{4pt}
\begin{minipage}[t]{0.17\textwidth}\vspace{0pt}\includegraphics[width=\linewidth]{alt_frames_v2/s39.png}\end{minipage}\hfill
\begin{minipage}[t]{0.80\textwidth}\vspace{0pt}\small
{\bfseries Scene 39: Where two mirrors agree}\hfill 2:50--2:54 (4\,s)\\[2pt]
\textbf{VO:} \emph{Not hidden. Just where two mirrors agree.}\\[2pt]
\textbf{On screen:} ``not hidden: where two mirrors agree''; reply to @iquilezles\\[2pt]
\textbf{Animation:} Final hold; loop.
\end{minipage}\par\vspace{4pt}\hrule\vspace{4pt}
\newpage
{\large\bfseries Full voiceover script} (174\,s, 434 words)\\[6pt]
1. [0:00--0:04] Inigo Quilez's video is lovely, and it opens many doors.\\
2. [0:04--0:08] I took one: circles as mirrors, the Möbius way.\\
3. [0:08--0:12] My aim: give his hidden points an algebra, and an invariant.\\
4. [0:12--0:17] His move: when the square root goes negative, go complex. Two orange dots.\\
5. [0:17--0:21] My move: a circle, a mirror-circle, and out comes another circle.\\
6. [0:21--0:25] A line is the simplest mirror: reflect a circle across it.\\
7. [0:25--0:29] A circle is a mirror too. That reflection is called inversion.\\
8. [0:29--0:33] Invert the line in the circle: out comes a small circle.\\
9. [0:33--0:38] Reflect the circle in the line instead: a different circle. Order matters.\\
10. [0:38--0:42] Every circle is just four numbers: a, b, c, d.\\
11. [0:42--0:46] And reflection is a single formula on those numbers.\\
12. [0:46--0:50] Circle in, circle out. Lines too. The family is closed.\\
13. [0:50--0:55] Rule one: a mirror fixes itself. Rule two: reflect twice, and you're back.\\
14. [0:55--1:00] Rule three: reflect a whole picture, and every relation in it comes along.\\
15. [1:00--1:04] These rules have a name: a kei, or involutory quandle.\\
16. [1:04--1:08] Not associative, though: move the brackets, and the circle changes.\\
17. [1:08--1:13] Points can't be mirrors, since everything would collapse. But they still get reflected.\\
18. [1:13--1:18] Now, the invariant: every mirror keeps the angle between any two circles.\\
19. [1:18--1:23] And that number decides real or hidden, so no mirror can flip it.\\
20. [1:23--1:27] What does a mirror leave alone? The points on it.\\
21. [1:27--1:31] And every circle that crosses it at right angles.\\
22. [1:31--1:35] Now take two mirrors: our circle and our line.\\
23. [1:35--1:39] Where do they agree? Same point in, same point out.\\
24. [1:39--1:44] When the line cuts, they agree at the crossings. Both leave them fixed.\\
25. [1:44--1:48] The circles both mirrors fix swirl around those two points.\\
26. [1:48--1:53] Lift the line, and the crossings hide. Do the mirrors still agree?\\
27. [1:53--1:58] They do, at two real points on the axis: the limiting points.\\
28. [1:58--2:03] Each mirror swaps them, top for bottom, and their heights multiply to one.\\
29. [2:03--2:07] Now the circles both mirrors fix all run through them.\\
30. [2:07--2:11] And Quilez's dots? The same two points, turned ninety degrees.\\
31. [2:11--2:16] Sweep the line: the points slide up, merge where it touches, then split.\\
32. [2:16--2:21] Two mirrors always agree at exactly two points, or one if they touch.\\
33. [2:21--2:27] Why two? Two reflections make one Möbius map. Its fixed points solve a quadratic.\\
34. [2:27--2:32] Over the complex numbers there are four: two crossings, two limiting points.\\
35. [2:32--2:36] But only two are ever real. The other two hide.\\
36. [2:36--2:41] Move everything by any mirror, and agreement points follow. That's the invariant answer.\\
37. [2:41--2:46] So the hidden pair isn't lost. It's the hidden half of the agreement.\\
38. [2:46--2:50] A circle, a mirror-circle, another circle: one closed algebra.\\
39. [2:50--2:54] Not hidden. Just where two mirrors agree.
\end{document}